%==============================================================================
%  APPENDIX - SOLID principles and design patterns
%  Application to the Zümm system: "before / after" justification
%==============================================================================

%==============================================================================
\chapter{Appendix -- SOLID principles and design patterns}
\label{chap:solid}
%==============================================================================
This appendix answers the question of the \emph{how} of software quality. It
documents the application of the five \textbf{SOLID} principles and a set of
\textbf{design patterns} to the \textbf{Zümm} system, setting each decision
against the initial design (\emph{before}) and the refactored design
(\emph{after}). The objective is a \textbf{multi-tier, scalable and loosely
coupled} architecture, compliant with the requirements of
chapter~\ref{chap:conception}.

\begin{encadre}
\textbf{Methodological sources:} the definitions adopted follow the presentation
of the SOLID principles from \emph{alexsoyes.com} and the pattern catalogue from
\emph{refactoring.guru}, adapted to the project's beekeeping domain.
\end{encadre}

%------------------------------------------------------------------------------
\section{The five SOLID principles}
%------------------------------------------------------------------------------
\begin{xltabular}{\textwidth}{>{\bfseries}p{2.4cm} p{6cm} X}
\toprule
\textbf{Principle} & \textbf{Before (initial design)} & \textbf{After (refactored)} \\
\midrule
\endhead
S: Single responsibility &
The \texttt{Ruche} entity carries the calculation method
\texttt{getZummHoneyActualQuantity()}: it mixes state (data) and business
logic. \texttt{Visite} groups data and report fields. &
The calculation is extracted into \texttt{ServiceMielImpl} (via the Composite)
and the report becomes \texttt{RapportVisite}. Each class now has only a single
reason to change. \\
O: Open / closed &
Alert detection and unit conversion would rely on \texttt{if / switch} on
\texttt{typeSeuil} and \texttt{unite}: adding a rule or a unit requires
modifying existing code. &
\texttt{RegleAlerte} (Strategy) and \texttt{ConvertisseurUnite}: a new rule or
unit = a new class, without touching the \texttt{GestionnaireAlertes}. Open to
extension, closed to modification. \\
L: Liskov substitution &
The use-case diagram generalises the supervising agent from the beekeeping
agent. Transposed as is into class inheritance, a restrictive subtype would
break substitutability. &
\texttt{Corps} and \texttt{Hausse} inherit from \texttt{Compartiment} while
fully respecting its contract (capacity 1..10, \texttt{getPoidsMiel}). The
agent's role goes through the \texttt{Role} enumeration + permissions
(composition) rather than a restrictive inheritance. \\
I: Interface segregation &
A single ``catch-all'' service (CRUD + visits + dashboards + honey API) would
force the third-party application to depend on unused methods. &
Targeted interfaces: \texttt{ServiceRequeteMiel} (the third-party API sees only
\texttt{getZummHoneyActualQuantity}), \texttt{ServiceVisite},
\texttt{ServiceTableauDeBord}. Each client depends only on its contract. \\
D: Dependency inversion &
The business layer (EJB) calls JDBC directly (concrete implementation): strong
coupling to the database and to \texttt{ResultSet}. &
The business layer depends on abstractions \texttt{RucheRepository} /
\texttt{VisiteRepository} (DAO pattern), implemented by
\texttt{RucheRepositoryJdbc}. Mock testing and JDBC \textrightarrow{} Spring
Data JPA migration facilitated (cf. appendix~\ref{chap:migration}). \\
\bottomrule
\end{xltabular}

%------------------------------------------------------------------------------
\section{Chosen design patterns}
%------------------------------------------------------------------------------
The following table lists the applied patterns, their category (creational,
structural, behavioural) and the layer concerned, with the \emph{before / after}
justification.

\begin{xltabular}{\textwidth}{>{\bfseries}p{2.7cm} p{5.9cm} X}
\toprule
\textbf{Pattern} & \textbf{Before} & \textbf{After} \\
\midrule
\endhead
Composite \newline {\small\itshape structural: business} &
The honey calculation manually traverses the frame collections with duplicated
summation code. &
\texttt{ElementRuche.getPoidsMiel()} is resolved recursively
(Ruche \textrightarrow{} Compartiment \textrightarrow{} Cadre): uniform
processing of the tree. \\
State \newline {\small\itshape behavioural: business} &
The hive's status would be tested by scattered \texttt{switch} (created, active,
harvesting, closed, \ldots). &
\texttt{EtatRuche} and its concrete states encapsulate the allowed transitions.
Direct alignment with the state-machine diagram. \\
Strategy \newline {\small\itshape behavioural: business} &
The alert rules, unit conversion and ROI would be hard-coded. &
Families of interchangeable algorithms (\texttt{RegleAlerte},
\texttt{ConvertisseurUnite}), injectable and testable in isolation. \\
Observer \newline {\small\itshape behavioural: business} &
The measurement code would call the dashboards directly (strong coupling). &
\texttt{GestionnaireAlertes} notifies the subscribed \texttt{Observateur}
(alerts dashboard, manager notification) when a threshold is crossed. \\
Command \newline {\small\itshape behavioural: control} &
The field operations (harvest, split, requeening) are direct calls, not
replayable offline. &
Each operation becomes a \texttt{CommandeVisite}. The
\texttt{FileCommandesHorsLigne} replays the commands on reconnection
(offline mode + deferred synchronisation). \\
Factory Method \newline {\small\itshape creational: business} &
The hives are created by scattered \texttt{new} with manual initialisation of
the compartments. &
\texttt{FabriqueRuche.creer(modele)} produces a preconfigured hive according to
the model (Dadant, Langstroth\ldots) with consistent body and capacities. \\
Builder \newline {\small\itshape creational: presentation} &
The visit report, with its many optional fields, would require a telescoping
constructor or scattered setters. &
\texttt{ConstructeurRapport} builds \texttt{RapportVisite} fluently
(findings, actions, photos, assessments). \\
Adapter \newline {\small\itshape structural: data} &
The heterogeneous sensors (different formats and units) would be handled case by
case. &
\texttt{AdaptateurCapteur} standardises each source into a \texttt{Mesure} with
configurable units (automated part). \\
Facade \newline {\small\itshape structural: control} &
The Servlet orchestrates all the business services directly. &
\texttt{FacadeApplication} offers a simplified entry point to the Servlets on
top of the subsystems. \\
Singleton \newline {\small\itshape creational: cross-cutting} &
The \texttt{ConfigZumm.ini} file and the i18n labels would be re-read in several
instances. &
\texttt{Configuration} and \texttt{GestionnaireI18n} guarantee a single shared
instance. \\
Proxy \newline {\small\itshape structural: web services} &
Access to the API would directly expose the business service. &
A \texttt{ProxySecurite} controls OAuth authentication and SSL / X.509
encryption before delegating to the service. \\
\bottomrule
\end{xltabular}

%------------------------------------------------------------------------------
\section{Refactored design view}
%------------------------------------------------------------------------------
Figure~\ref{fig:solid} summarises the main patterns in a design class diagram:
the Composite (hive composition and honey calculation), the State (life cycle),
the segregated services (ISP), abstraction-based persistence (DIP / Repository),
the Strategy + Observer pair (alerts) and the Factory / Builder / Command trio.

\begin{figure}[htbp]
\centering
\includegraphics[width=\linewidth]{diagramme_classes_solid}
\caption{Refactored class diagram: application of SOLID and the design patterns
to the Zümm system.}
\label{fig:solid}
\end{figure}

\begin{encadre}
\textbf{Correspondence with the layers:} the \texttt{Repository} belong to the
data-access layer (JDBC), the \texttt{Service} and the domain (Composite, State,
Strategy, Observer, Command) to the business layer (EJB), the \texttt{Facade}
and the \texttt{Proxy} to the control layer (Servlets), the report
\texttt{Builder} to the presentation (JSP). The loosely coupled layered
architecture of chapter~\ref{chap:conception} is thus reinforced, without being
modified.
\end{encadre}

%------------------------------------------------------------------------------
\section{Dynamic illustration: Observer and Command patterns}
%------------------------------------------------------------------------------
Two sequence diagrams detail the behaviour of the system's most structuring
patterns.

\subsection{Observer (+ Strategy): triggering an alert}
Figure~\ref{fig:seq-observer} shows how a measurement received by the
\texttt{GestionnaireAlertes} (subject) is evaluated by an interchangeable rule
(\emph{Strategy}) then, in case of a breach, propagated to the subscribed
observers (alerts dashboard, manager notification). Adding a rule or an observer
requires no modification of the manager (open/closed principle).

\begin{figure}[htbp]
\centering
\includegraphics[width=\linewidth]{diagramme_sequence_observer}
\caption{Sequence of the Observer (+ Strategy) pattern: alert-triggering flow.}
\label{fig:seq-observer}
\end{figure}

\subsection{Command: offline mode and deferred synchronisation}
Figure~\ref{fig:seq-command} illustrates the entry of a field operation without
a network: the operation is encapsulated in a \texttt{CommandeVisite} pushed
onto the \texttt{FileCommandesHorsLigne}. When the network returns, the queue
replays (\texttt{executer}) each command, or compensates it (\texttt{annuler})
in case of conflict, achieving the deferred synchronisation required by the
offline mode.

\begin{figure}[htbp]
\centering
\includegraphics[width=\linewidth]{diagramme_sequence_command}
\caption{Sequence of the Command pattern: offline mode and deferred synchronisation.}
\label{fig:seq-command}
\end{figure}

\subsection{State: hive life cycle}
Figure~\ref{fig:seq-state} shows the Ruche (context) delegating each operation to
its current state. The state decides the transition (activation, harvesting,
return to the active state) or refuses an operation forbidden in the current
state, here a populating requested on an already active hive. The behaviour thus
depends on the state without any \texttt{switch} on the status, consistently with
the state-machine diagram (figure~\ref{fig:etat}).

\begin{figure}[htbp]
\centering
\includegraphics[width=0.9\linewidth]{diagramme_sequence_state}
\caption{Sequence of the State pattern: delegation of the transitions of the hive life cycle.}
\label{fig:seq-state}
\end{figure}

\subsection{Composite: honey-quantity calculation}
Figure~\ref{fig:seq-composite} details the recursive resolution of
\texttt{getPoidsMiel()}: the client (\texttt{ServiceMielImpl}) queries the hive,
which propagates the call to each compartment, then to each frame (leaf). Only
frames of type \texttt{MIEL} contribute, and the tare is subtracted at the hive
level. The client ignores the internal structure. This calculation implements
the \texttt{getZummHoneyActualQuantity} method exposed by the web service.

\begin{figure}[htbp]
\centering
\includegraphics[width=0.85\linewidth]{diagramme_sequence_composite}
\caption{Sequence of the Composite pattern: recursive calculation of a hive's honey quantity.}
\label{fig:seq-composite}
\end{figure}

\subsection{Strategy: conversion of heterogeneous units}
Figure~\ref{fig:seq-strategy} shows the normalisation of a measurement expressed
in any unit: the \texttt{RegistreConvertisseurs} (context) selects the strategy
corresponding to the unit (\texttt{ConvertisseurLbVersKg}), which brings the
value back to the reference unit before comparison with the threshold. A new
unit translates into a new strategy, without modifying the context.

\begin{figure}[htbp]
\centering
\includegraphics[width=0.85\linewidth]{diagramme_sequence_strategy}
\caption{Sequence of the Strategy pattern: conversion of heterogeneous units to the reference unit.}
\label{fig:seq-strategy}
\end{figure}

\subsection{Adapter: ingestion of heterogeneous sensors}
Figure~\ref{fig:seq-adapter} illustrates the integration of a market sensor: the
\texttt{AdaptateurBroodMinder} translates the sensor's proprietary API into a
standardised \texttt{Mesure} (converted units, timestamp, geolocation),
exposing the uniform \texttt{Capteur} interface expected by the business layer.
The heterogeneous sensors thus become interchangeable.

\begin{figure}[htbp]
\centering
\includegraphics[width=\linewidth]{diagramme_sequence_adapter}
\caption{Sequence of the Adapter pattern: ingestion of heterogeneous sensors (automated part).}
\label{fig:seq-adapter}
\end{figure}

%------------------------------------------------------------------------------
\section{Distribution of the patterns per layer}
%------------------------------------------------------------------------------
Figure~\ref{fig:packages-patterns} summarises the location of each pattern in the
multi-tier architecture. The \emph{Builder} sits in presentation. The
\emph{Facade}, the \emph{Proxy} and the \emph{Command} queue belong to control.
The domain (\emph{Composite}, \emph{State}, \emph{Observer}, \emph{Strategy},
\emph{Factory Method}) and the segregated services (\emph{ISP}) are placed in
business. The \emph{Repository} (\emph{DAO / DIP}) and the \emph{Adapter} handle
data access. Finally, the configuration and internationalisation \emph{Singleton}
as well as the unit-conversion \emph{Strategy} occupy the cross-cutting layer.
This view confirms that the patterns reinforce the loosely coupled layered
architecture without altering its structure.

\begin{figure}[htbp]
\centering
\includegraphics[width=\linewidth]{diagramme_packages_patterns}
\caption{Distribution of the design patterns per layer of the multi-tier architecture.}
\label{fig:packages-patterns}
\end{figure}
